% GENERATED FILE. Edit canonical lessons, then run npm run generate:books.
% canonical-source-hash: fnv1a64:929771dbb8260ffd
% canonical-lessons: AR-C09-aswad-abyad, AR-C09-ahmar-azraq

\chapter{Four Core Colours}
\label{ch:ar-core-colours}

This chapter is generated from the canonical micro-lessons used by Language Ladder.
Each section stays independently resumable and preserves the authored prerequisite order.

\section[\ar{أسود} \ar{أبيض}]{\ar{أسود} / \ar{أبيض} (aswad / abyad) — black and white, a whole new gender pattern}

\label{lesson:AR-C09-aswad-abyad}



\begin{quote}
\textbf{Warm-up.} {[}PAUSE 2s{]} You've already met a gendered adjective — \textbf{\ar{آسف}/\ar{آسفة}} ("sorry"), which just added a feminine ending. Color words go further: they reshape the \textbf{whole word}, front and back, for gender.
\end{quote}

\begin{grammarlens}[title={Grammar Lens — The pattern: masculine \ar{أ}‑‑‑ / feminine ‑‑‑\ar{اء}}]
\begin{itemize}
\raggedright
  \item \textbf{\ar{أسود}} (\emph{aswad}, masculine) / \textbf{\ar{سوداء}} (\emph{sawdāʾ}, feminine) — "\textbf{black}"
  \item \textbf{\ar{أبيض}} (\emph{abyad}, masculine) / \textbf{\ar{بيضاء}} (\emph{bayḍāʾ}, feminine) — "\textbf{white}"
\end{itemize}

Notice the shape: the masculine form \textbf{starts} with \textbf{\ar{أ}} (a hamza-alif), and the feminine form \textbf{drops} that opening letter but \textbf{adds} \textbf{‑\ar{اء}} (\emph{-āʾ}) at the end instead. This isn't the simple "add one letter" pattern \emph{āsif/ āsifa} used — it's a \textbf{dedicated template} Arabic reserves specifically for \textbf{colors and physical characteristics} (this same shape appears again for red and blue in the next lesson).
\end{grammarlens}

\subsection*{Guided Practice}
{[}PAUSE 1s{]}

\begin{itemize}
\raggedright
  \item {[}YOU SAY: masculine — "aswad, abyad" — black, white{]}
  \item {[}YOU SAY: feminine — "sawdāʾ, bayḍāʾ" — black, white{]}
  \item {[}YOU SAY: notice the shape — \ar{أ}‑‑‑ drops, ‑‑‑\ar{اء} appears{]}
\end{itemize}

\begin{tcolorbox}[breakable,colback=teal!4,colframe=teal!35!black,title={Wrap-up recall}]
{[}PAUSE 3s{]} What's the masculine/feminine pair for black? (\textbf{\ar{أسود} aswad / \ar{سوداء} sawdāʾ}.) For white? (\textbf{\ar{أبيض} abyad / \ar{بيضاء} bayḍāʾ}.) How is this gender-shape different from \emph{āsif/āsifa}? (\textbf{It reshapes the whole word} — losing the opening \ar{أ} and gaining a closing ‑\ar{اء} — a template reserved for colors and physical traits, not just adding a simple ending.)
\end{tcolorbox}
\section[\ar{أحمر} \ar{أزرق}]{\ar{أحمر} / \ar{أزرق} (ahmar / azraq) — red and blue, the same template again}

\label{lesson:AR-C09-ahmar-azraq}



\begin{quote}
\textbf{Warm-up.} {[}PAUSE 2s{]} Two more colors — and the exact same shape you just learned.
\end{quote}

\begin{grammarlens}[title={Grammar Lens — The pattern holds}]
\begin{itemize}
\raggedright
  \item \textbf{\ar{أحمر}} (\emph{aḥmar}, masculine) / \textbf{\ar{حمراء}} (\emph{ḥamrāʾ}, feminine) — "\textbf{red}"
  \item \textbf{\ar{أزرق}} (\emph{azraq}, masculine) / \textbf{\ar{زرقاء}} (\emph{zarqāʾ}, feminine) — "\textbf{blue}"
\end{itemize}

Same recipe every time: masculine opens with \textbf{\ar{أ}}, feminine drops that opening and closes with \textbf{‑\ar{اء}} instead. Once you've internalized this template from black and white, the whole family of \textbf{classic, old} Arabic colors slots into it (yellow \textbf{\ar{أصفر}/\ar{صفراء}}, green \textbf{\ar{أخضر}/\ar{خضراء}}, too) — one pattern, several colors. (Newer colors like brown, gray, purple, orange, and pink use a \textbf{different} pattern entirely, borrowed from nouns — so "every color" would overstate it; this template belongs specifically to Arabic's oldest, core color set.)
\end{grammarlens}

\subsection*{Guided Practice}
{[}PAUSE 1s{]}

\begin{itemize}
\raggedright
  \item {[}YOU SAY: masculine — "aḥmar, azraq" — red, blue{]}
  \item {[}YOU SAY: feminine — "ḥamrāʾ, zarqāʾ" — red, blue{]}
  \item {[}YOU SAY: all four masculine forms in a row — "aswad, abyad, aḥmar, azraq"{]}
\end{itemize}

\begin{tcolorbox}[breakable,colback=teal!4,colframe=teal!35!black,title={Wrap-up recall}]
{[}PAUSE 3s{]} What's the masculine/feminine pair for red? (\textbf{\ar{أحمر} aḥmar / \ar{حمراء} ḥamrāʾ}.) For blue? (\textbf{\ar{أزرق} azraq / \ar{زرقاء} zarqāʾ}.) Does every Arabic color use this same template? (\textbf{No} — Arabic's classic, oldest colors (black, white, red, blue, yellow, green) share \textbf{one} template, masc. \ar{أ}‑‑‑ / fem. ‑‑‑\ar{اء}; newer colors like brown or purple use a different, noun-based pattern instead.)
\end{tcolorbox}
